\documentclass[10pt,twocolumn]{article}
\usepackage{listings}
\usepackage{xcolor}
\usepackage[margin=1cm]{geometry}
\usepackage{longtable,booktabs,array}
\newcounter{none}
\providecommand{\passthrough}[1]{#1}
\providecommand{\tightlist}{\setlength{\itemsep}{0pt}\setlength{\parskip}{0pt}}
\input{fun}

\begin{document}

\section{NAME}\label{name}

\textbf{fun} --- transpile and run a \emph{fun} source file (a tiny
Lua-targeted source language)

\section{SYNOPSIS}\label{synopsis}

\textbf{fun} \emph{FILE} {[}\emph{ARGS}\ldots{]}

\section{DESCRIPTION}\label{description}

\emph{fun} is a Lua-targeted source language whose transpiler is a
single short Lua module (\passthrough{\lstinline!fun.lua!}). Source
files use the \passthrough{\lstinline!.fun!} extension. The runner shim
\passthrough{\lstinline!fun!} reads a \passthrough{\lstinline!.fun!}
file, transpiles it to Lua via \passthrough{\lstinline!fun.lua!},
executes the result with the remaining arguments visible in Lua's
\passthrough{\lstinline!arg!} table, and forwards stdin/stdout/exit
code.

By the time the source reaches Lua, it \emph{is} Lua. \emph{fun} adds a
small set of line-level rewrites that compile away --- no runtime layer.

\section{QUICK START}\label{quick-start}

\begin{lstlisting}[language=bash]
./fun foo.fun arg1 arg2
\end{lstlisting}

Inside \passthrough{\lstinline!foo.fun!}:

\begin{lstlisting}[language=Fun]
#!/usr/bin/env ./fun
-- vim: ft=fun

let greet, sq

greet = fun(name) ! "hello, "..name end
sq    = fun(x)    ! x*x end

print(greet("world"))
print(sq(5))
\end{lstlisting}

\section{LANGUAGE}\label{language}

\subsection{Keyword sugar}\label{keyword-sugar}

{\def\LTcaptype{none} % do not increment counter
\begin{longtable}[]{@{}lll@{}}
\toprule\noalign{}
\emph{fun} & Lua & Notes \\
\midrule\noalign{}
\endhead
\bottomrule\noalign{}
\endlastfoot
\passthrough{\lstinline!fun!} & \passthrough{\lstinline!function!} &
word-boundary substitution \\
\passthrough{\lstinline!let!} & \passthrough{\lstinline!local!} &
word-boundary substitution \\
\passthrough{\lstinline"!"} & \passthrough{\lstinline!return!} & sigil;
emits \passthrough{\lstinline!return!} + RHS \\
\end{longtable}
}

Substitutions are line-level. Strings (\passthrough{\lstinline!"..."!},
\passthrough{\lstinline!'...'!}, \passthrough{\lstinline![[...]]!}) and
trailing \passthrough{\lstinline!--!} comments are hidden before
substitution and restored after, so \passthrough{\lstinline'"!"'} and
other token-like contents survive unchanged.

\subsection{Control flow}\label{control-flow}

\passthrough{\lstinline!if!} and \passthrough{\lstinline!elseif!}
\textbf{require parentheses} around the condition; the transpiler
injects \passthrough{\lstinline!then!}. \passthrough{\lstinline!end!}
and \passthrough{\lstinline!do!} are written literally.

\begin{lstlisting}[language=Fun]
if (x > 0)
  print("pos")
elseif (x < 0)
  print("neg")
else
  print("zero")
end
\end{lstlisting}

\passthrough{\lstinline!for!}, \passthrough{\lstinline!while!},
\passthrough{\lstinline!repeat ... until!} use plain Lua syntax
(\passthrough{\lstinline!do!}/\passthrough{\lstinline!end!} explicit, no
sugar).

\subsection{Function definitions}\label{function-definitions}

\passthrough{\lstinline!fun NAME(args) ... end!} works like Lua's
\passthrough{\lstinline!function!}:

\begin{lstlisting}[language=Fun]
fun sq(n) ! n*n end
fun NUM.mid(i) ! i.mu end
\end{lstlisting}

Anonymous: \passthrough{\lstinline!fun(args) ... end!}.

Named-as-local: combine with \passthrough{\lstinline!let!}:

\begin{lstlisting}[language=Fun]
let helper = fun(x) ! x+1 end
\end{lstlisting}

A trailing \passthrough{\lstinline!=!} on a
\passthrough{\lstinline!fun(args)!} line is silently stripped (so an old
\passthrough{\lstinline!name = fun(args) =!} style still parses).

\subsection{\texorpdfstring{Conditional init:
\texttt{?=}}{Conditional init: ?=}}\label{conditional-init}

\begin{lstlisting}[language=Fun]
x ?= 42      -- if x == nil then x = 42 end
\end{lstlisting}

Statement-only; must start the line (optional indent allowed).

\subsection{Compound assignment}\label{compound-assignment}

\begin{lstlisting}[language=Fun]
x += 1       -- x = x + 1
y -= 2       -- y = y - 2
z *= 3       -- etc.
w /= 4
\end{lstlisting}

Operators: \passthrough{\lstinline!+= -= *= /=!}. \textbf{LHS must be a
single word} (\passthrough{\lstinline!\%w+!}) and \textbf{RHS must be
one whitespace-separated token} (\passthrough{\lstinline!\%S+!}).

\subsection{Line continuation}\label{line-continuation}

A trailing \passthrough{\lstinline!\\!} joins the next line; the
joined-from physical line is replaced with a blank to \textbf{preserve
line numbers} for error reporting.

\begin{lstlisting}[language=Fun]
let result = some_function(arg1, arg2, \
                           arg3, arg4)
\end{lstlisting}

This is the \textbf{only} way to span constructs across physical lines
for features whose regex is per-line --- see \emph{one-line constraints}
below.

\subsection{Comprehensions}\label{comprehensions}

\begin{lstlisting}[language=Fun]
let squares = [x*x for x in xs]
let evens   = [x   for x in xs if x % 2 == 0]
\end{lstlisting}

Compiles to an inline IIFE (immediately invoked function expression). If
\passthrough{\lstinline!iter!} contains no \passthrough{\lstinline!(!},
the transpiler wraps:

\begin{itemize}
\tightlist
\item
  single loop var → \passthrough{\lstinline!for \_,v in ipairs(iter)!}
\item
  comma in loop var → \passthrough{\lstinline!for k,v in pairs(iter)!}
\end{itemize}

If \passthrough{\lstinline!iter!} already contains
\passthrough{\lstinline!(!}, it passes through verbatim.

\subsection{Strings and comments}\label{strings-and-comments}

\begin{itemize}
\tightlist
\item
  \passthrough{\lstinline!--!} starts a single-line comment.
\item
  \passthrough{\lstinline!--[[ ... ]]!} is a long comment.
\item
  \passthrough{\lstinline![[ ... ]]!} is a long string (multi-line,
  \textbf{no interpolation}).
\item
  \passthrough{\lstinline!"..."!}, \passthrough{\lstinline!'...'!} are
  regular strings (\textbf{no interpolation}).
\item
  All string contents protected from substitution.
\end{itemize}

\subsection{Shebang}\label{shebang}

\passthrough{\lstinline"\#!/usr/bin/env ./fun"} on line 1 makes the file
directly executable; the transpiler converts the shebang to a Lua
comment before loading.

\section{OPTIONS}\label{options}

The \passthrough{\lstinline!fun!} runner takes no flags of its own. The
first positional argument is the \passthrough{\lstinline!.fun!} file;
remaining arguments are forwarded to the script via
\passthrough{\lstinline!arg[1]!}, \passthrough{\lstinline!arg[2]!},
\ldots{} \passthrough{\lstinline!arg[0]!} is set to the script path.

\section{FILES}\label{files}

\begin{itemize}
\tightlist
\item
  \passthrough{\lstinline!fun.lua!} --- the transpiler (a single Lua
  module).
\item
  \passthrough{\lstinline!fun!} --- the runner shim
  (\passthrough{\lstinline"\#!/usr/bin/env lua"}).
\item
  \passthrough{\lstinline!*.fun!} --- \emph{fun} source files.
\item
  \passthrough{\lstinline!etc/funlexer.py!} --- Pygments lexer
  (\passthrough{\lstinline!FunLexer!}).
\item
  \passthrough{\lstinline!etc/funpdf.sh!} ---
  \passthrough{\lstinline!.fun!} → PDF via pygmentize + pandoc +
  pdflatex.
\item
  \passthrough{\lstinline!etc/syntax/fun.vim!} --- vim syntax (loaded
  via \passthrough{\lstinline!etc/nvimfun.lua!}).
\item
  \passthrough{\lstinline!etc/nvimfun.lua!},
  \passthrough{\lstinline!etc/nviminit.lua!} --- neovim filetype plugin.
\end{itemize}

\section{DIAGNOSTICS}\label{diagnostics}

When the transpiled Lua fails to load, \passthrough{\lstinline!fun!}
writes a 5-line window around the offending line (\textbf{transpiled}
Lua, not source) to stderr, then errors out. Line numbers refer to
\textbf{transpiled} Lua, but the blank-pad trick on
\passthrough{\lstinline!\\!} joins keeps source and transpiled line
numbers in lockstep, so the position usually points to the right
physical source line.

\section{ONE-LINE CONSTRAINTS}\label{one-line-constraints}

The transpiler runs substitutions \textbf{per physical line}
(\passthrough{\lstinline!src:gsub "[\^\\n]+"!}). Several constructs must
therefore fit on one line. Workaround: use \passthrough{\lstinline!\\!}
continuation.

\subsection{\texorpdfstring{\texttt{if\ (...)} / \texttt{elseif\ (...)}
--- one
line}{if (...) / elseif (...) --- one line}}\label{if-...-elseif-...-one-line}

The injection regex matches \passthrough{\lstinline!if!} (or
\passthrough{\lstinline!elseif!}) followed by a balanced-paren
expression on the same line. Multi-line conditions need
\passthrough{\lstinline!\\!}:

\begin{lstlisting}[language=Fun]
-- ✗ broken — closing paren on next physical line
if (very_long_condition_part1 and
    very_long_condition_part2)
  ...

-- ✓ join with backslash
if (very_long_condition_part1 and \
    very_long_condition_part2)
  ...
\end{lstlisting}

\subsection{List comprehensions --- one
line}\label{list-comprehensions-one-line}

\passthrough{\lstinline![expr for var in iter]!} (and the
\passthrough{\lstinline!if!}-guarded form) must fit on one line. The
regex stops at the first matching \passthrough{\lstinline!]!}.

\begin{lstlisting}[language=Fun]
-- ✗ broken — comprehension spans lines
let xs = [transform(x)
          for x in source]

-- ✓ join with backslash
let xs = [transform(x) \
          for x in source]

-- ✓ or extract the iterator first
let it = source
let xs = [transform(x) for x in it]
\end{lstlisting}

\subsection{\texorpdfstring{\texttt{?=} --- one line,
statement-leading}{?= --- one line, statement-leading}}\label{one-line-statement-leading}

Pattern: \passthrough{\lstinline!\^(\%s*)(\%w+)\%s*?=\%s*([\^;]+)!}.
Indented allowed; the RHS runs up to \passthrough{\lstinline!;!} or
end-of-line. Cannot follow other content on the same line.

\subsection{\texorpdfstring{Compound
\texttt{+=}/\texttt{-=}/\texttt{*=}/\texttt{/=} --- one word
LHS}{Compound +=/-=/*=//= --- one word LHS}}\label{compound---one-word-lhs}

\passthrough{\lstinline!(\%w+)\%s*([\%+\%-\%*/])=\%s+(\%S+)!}. Both
sides constrained:

\begin{itemize}
\tightlist
\item
  LHS is a single \passthrough{\lstinline!\%w+!} --- \textbf{no fields
  or indexing}. \passthrough{\lstinline!obj.field += 1!} and
  \passthrough{\lstinline!t[i] += 1!} do \textbf{not} transform.
\item
  RHS is \textbf{one} whitespace-separated token
  (\passthrough{\lstinline!\%S+!}). \passthrough{\lstinline!x += a + b!}
  captures only \passthrough{\lstinline!a!}; the rest leaks.
\end{itemize}

Workaround: write the expansion by hand or wrap RHS in parens with no
internal whitespace --- \passthrough{\lstinline!x += (a+b)!} works
because \passthrough{\lstinline!(a+b)!} is one token.

\subsection{\texorpdfstring{Comprehension iterator with \texttt{{]}}
inside}{Comprehension iterator with {]} inside}}\label{comprehension-iterator-with-inside}

The regex stops at the \textbf{first} \passthrough{\lstinline!]!}. If
\passthrough{\lstinline!iter!} contains a literal
\passthrough{\lstinline!]!} (e.g.~inside a Lua character class), the
regex closes early:

\begin{lstlisting}[language=Fun]
-- ✗ broken — `]` inside "[^,]+" closes the comprehension
let t = [l.thing(x) for x in s:gmatch"[^,]+"]
\end{lstlisting}

Workaround: bind the iterator first.

\begin{lstlisting}[language=Fun]
let it = s:gmatch"[^,]+"
let t  = [l.thing(x) for x in it]
\end{lstlisting}

\section{OTHER QUIRKS}\label{other-quirks}

\subsection{\texorpdfstring{\texttt{!} substitutes everywhere outside
strings/comments}{! substitutes everywhere outside strings/comments}}\label{substitutes-everywhere-outside-stringscomments}

\passthrough{\lstinline"!"} is a literal global substitution to
\passthrough{\lstinline!return!} (with trailing space). Lua identifiers
can't contain \passthrough{\lstinline"!"}, so accidental damage is rare,
\textbf{but}:

\begin{itemize}
\tightlist
\item
  \passthrough{\lstinline"!="} becomes
  \passthrough{\lstinline!return =!} --- broken. Use Lua's
  \passthrough{\lstinline!\~=!} for not-equal.
\item
  \passthrough{\lstinline"not !x"} becomes
  \passthrough{\lstinline!not return x!} --- broken. Just write
  \passthrough{\lstinline!not x!}.
\item
  \passthrough{\lstinline"if (x != y)"} is \textbf{wrong}; write
  \passthrough{\lstinline!if (x \~= y)!}.
\end{itemize}

\subsection{\texorpdfstring{\texttt{fun} and \texttt{let} use word
boundaries}{fun and let use word boundaries}}\label{fun-and-let-use-word-boundaries}

Frontier patterns \passthrough{\lstinline!\%f[\%w\_]fun\%f[\%W]!} and
\passthrough{\lstinline!\%f[\%w\_]let\%f[\%W]!}. So:

\begin{itemize}
\tightlist
\item
  \passthrough{\lstinline!letx!} is \textbf{not} touched (no boundary).
\item
  \passthrough{\lstinline!myfun!} is \textbf{not} touched.
\item
  \passthrough{\lstinline!funny!} is \textbf{not} touched
  (\passthrough{\lstinline!y!} extends word).
\item
  \passthrough{\lstinline!fun(!} and \passthrough{\lstinline!fun!} and
  \passthrough{\lstinline!let!} are transformed.
\end{itemize}

\subsection{\texorpdfstring{String hiding does not handle
\texttt{\textbackslash{}"}-escaped
quotes}{String hiding does not handle \textbackslash"-escaped quotes}}\label{string-hiding-does-not-handle--escaped-quotes}

The hiding regex is \passthrough{\lstinline!'"[\^"]*"'!}. Strings
containing escaped quotes break the hiding:

\begin{lstlisting}[language=Fun]
let bad = "she said \"hi\""    -- hiding misparses this
\end{lstlisting}

Workaround: use single quotes, \passthrough{\lstinline![[...]]!}, or
split into pieces with Lua's \passthrough{\lstinline!..!} concat.

\subsection{Method dispatch on plain
tables}\label{method-dispatch-on-plain-tables}

Plain Lua tables have no methods. \passthrough{\lstinline!xs:remove(n)!}
is \textbf{not} sugar for \passthrough{\lstinline!table.remove(xs, n)!}
--- Lua looks for a \passthrough{\lstinline!remove!} field and errors.
Use \passthrough{\lstinline!table.remove(xs, n)!} or alias at top:
\passthrough{\lstinline!let tremove = table.remove!}.

\subsection{No string interpolation
anywhere}\label{no-string-interpolation-anywhere}

Neither \passthrough{\lstinline!"..."!},
\passthrough{\lstinline!'...'!}, nor \passthrough{\lstinline![[...]]!}
interpolate \passthrough{\lstinline!\{var\}!}. Build strings with Lua's
\passthrough{\lstinline!..!} concat.

\subsection{Long strings span multiple lines
safely}\label{long-strings-span-multiple-lines-safely}

\passthrough{\lstinline![[...]]!} blocks are extracted from the whole
source before line-by- line processing, so multi-line long strings (and
blank lines inside them) work fine.

\subsection{Comments are passed through
verbatim}\label{comments-are-passed-through-verbatim}

A trailing \passthrough{\lstinline!--!} comment is split off before
substitution and re- appended after, so anything inside a comment
survives unchanged.

\section{EXAMPLE}\label{example}

\begin{lstlisting}[language=Fun]
#!/usr/bin/env ./fun
-- vim: ft=fun

let Num, add

Num = fun(s,n) ! {name=s, at=n, n=0, mu=0, m2=0, sd=0} end

add = fun(c, v)
  c.n += 1
  let d = v - c.mu
  c.mu += d / c.n
  c.m2 += d * (v - c.mu)
  c.sd  = c.n < 2 and 0 or (c.m2 / (c.n - 1)) ^ 0.5
  ! c
end

-- comprehension: build cumulative stats
let xs = [add(Num("x",0), v) for v in {1,2,3,4,5}]
print(xs[#xs].mu, xs[#xs].sd)
\end{lstlisting}

\section{SEE ALSO}\label{see-also}

\passthrough{\lstinline!lua(1)!}, \passthrough{\lstinline!pairs(3)!},
\passthrough{\lstinline!ipairs(3)!},
\passthrough{\lstinline!pygmentize(1)!},
\passthrough{\lstinline!pandoc(1)!}

\section{AUTHORS}\label{authors}

Tim Menzies \href{mailto:timm@ieee.org}{\nolinkurl{timm@ieee.org}}, with
iterative collaboration.

\end{document}
